\chapter{Introducción}
\label{cap:introduccion}

\section{Envejecimiento y protección económica en México}
\label{sec:problema-proteccion}

El envejecimiento demográfico hace necesario examinar tanto el acceso a recursos
durante la vejez como las diferencias en su cuantía y continuidad. La Encuesta
Nacional de la Dinámica Demográfica (ENADID) 2023 estimó que la proporción
de personas de 60 años o más en México aumentó
de 12,3\% en 2018 a 14,7\% en 2023 \citep[p.~2]{inegi2024enadid}.
El dato sitúa el problema en una transformación de la estructura por edad;
no implica que todas las personas mayores tengan las mismas necesidades,
recursos o trayectorias laborales. La referencia de 2023 ofrece contexto
demográfico, mientras que el periodo empírico de esta investigación termina
en la ronda nominal de 2021.

En la vejez pueden coexistir ingresos laborales, pensiones, transferencias
públicas, ayuda familiar y utilización de activos acumulados. Estas fuentes
no son intercambiables. Una prestación periódica representa un flujo; el
patrimonio es un stock y puede estar compartido o tener liquidez limitada.
Del mismo modo, que una persona resida con alguien pensionado no acredita
que tenga un derecho propio o que disponga de una proporción determinada de
ese ingreso. La investigación distingue por ello recepción personal,
recepción del cónyuge y recursos económicos totales, estos últimos fuera del
alcance de los indicadores monetarios aquí analizados.

El sistema mexicano reúne instituciones, contingencias y reglas de acceso
heterogéneas. Las pensiones pueden corresponder a retiro, invalidez,
accidente o derechos de supervivencia; los apoyos públicos pueden incluir
prestaciones de vejez, otras transferencias monetarias y recursos en especie.
La expansión de los programas federales y locales durante 2001--2021 modificó
edades y ámbitos de elegibilidad, importes y posibilidades de coexistencia
con otras prestaciones \citep{setentamas2007,setentamas2012,pam2013,
pam2014,bienestar2019,bienestar2021b}. Esa historia motiva estudiar conjuntamente
quién recibe y cuánto recibe, sin atribuir a una regla la variación observada
en una encuesta por el solo hecho de coincidir en el tiempo.

\section{Problema de investigación}

Ampliar la recepción de alguna prestación puede reducir diferencias de acceso
sin igualar los montos entre sus receptores. Si quienes se incorporan tienen
importes menores que quienes ya declaraban recepción, la distribución del
conjunto cambia aun cuando los continuos no experimenten una pérdida
equivalente. A su vez, un incremento nominal puede resultar insuficiente
para mantener el poder adquisitivo. Son procesos observacionalmente distintos:
cambio de cobertura, cambio de composición, variación de la distribución
monetaria y cambio de la misma persona.

El problema requiere además separar protección propia de contexto familiar.
La pensión del cónyuge puede constituir un recurso relevante en el hogar,
pero no se transforma por ello en pensión del entrevistado. Tampoco la
negativa de una pregunta pensionaria equivale a carecer de cualquier apoyo:
la recepción puede aparecer en otra sección, permanecer desconocida o no ser
identificable bajo la definición utilizada. La ausencia de una vía propia
K/A es, por tanto, una condición de información y recepción definidas,
no un diagnóstico de desprotección absoluta.

Una comparación longitudinal añade otra dificultad. Los receptores
observados en dos entrevistas son una selección de quienes estaban al inicio.
Hay fallecimientos registrados, ausencias cuyo estado vital no se conoce,
entrevistas por sustituto y respuestas monetarias incompletas. Estos destinos
no deben reunirse bajo una única categoría de faltantes. Un cambio monetario
solo calculable entre sobrevivientes reentrevistados no describe automáticamente
a toda la cohorte, y una prestación posterior al fallecimiento no constituye
un valor faltante que deba imputarse.

La tesis aborda estas distinciones mediante el Estudio Nacional de Salud y
Envejecimiento en México (ENASEM), para 2001--2021 \citep{wong2017}. Emplea las seis
rondas para describir el alcance de las fuentes disponibles y restringe las
comparaciones temporales a ventanas con comparabilidad operacional condicionada.
El mayor detalle monetario longitudinal se concentra en K 2012--2015 y
A 2015--2018. El periodo general del estudio no implica que exista una serie
homogénea de protección económica durante veinte años.

\section{Pregunta de investigación}

La investigación conserva la siguiente pregunta:

\begin{quote}
¿Qué tanto es consistente la expansión de las pensiones y los apoyos públicos
con una reducción de las brechas de cobertura y de monto de la protección
económica de los adultos mayores en México (ENASEM 2001--2021), qué grupos siguen
sin vía propia de protección y cómo cambia esta lectura al considerar a
quienes mueren o dejan de ser entrevistados?
\end{quote}

La expresión \emph{consistente con} delimita una comparación empírica:
se examina si acceso, distribución monetaria y trayectorias observadas presentan
patrones compatibles entre sí. No se identifica el efecto de una reforma.
Las referencias a México indican el contexto institucional y la encuesta;
los resultados propios describen muestras analíticas sin ponderadores o
diseño nacional validados para estos estimandos.

\section{Justificación económica, estadística y actuarial}

La justificación económica reside en distinguir acceso de recursos recibidos.
Una tasa de recepción informa sobre una vía de protección; una mediana
condicionada a recepción caracteriza un nivel central entre beneficiarios.
Ninguna resume por sí sola la suficiencia de ingreso de una persona o un hogar.
Compararlas permite reconocer perfiles con mayor acceso a alguna prestación
y, simultáneamente, importes relativamente bajos.

La justificación estadística es atender la población de cada indicador.
Los denominadores cambian cuando se requiere respuesta conocida, un importe
válido, información geográfica para el umbral o dos mediciones de la misma
persona. Las medianas son pertinentes ante distribuciones monetarias
asimétricas, pero su coincidencia no demuestra igualdad distributiva. Los
empates, la información incompleta y la composición de los grupos deben
acompañar cualquier intervalo de incertidumbre.

Desde la Actuaría interesa cómo los recursos se observan durante el
envejecimiento y qué población queda representada cuando se condiciona al
seguimiento. La muerte trunca la trayectoria individual; una ausencia puede
ser temporal; un cambio de informante modifica la observabilidad sin
demostrar que cesó la prestación. Esta perspectiva fundamenta sensibilidades
de selección y continuidad, manteniendo separados los procesos. La actualización
por inflación conserva poder adquisitivo comparable; no representa el valor
presente actuarial de una corriente vitalicia.

La contribución de Ciencia de Datos consiste en una cadena reproducible
de identificación, titularidad, eventos, estados desconocidos y agregados.
Esa construcción permite relacionar la evidencia transversal y longitudinal
sin duplicar beneficios ni convertir la información del hogar en ingreso
individual. La reproducibilidad sirve al argumento sustantivo y no sustituye
la validación conceptual de las preguntas.

\section{Objetivos}

\subsection{Objetivo general}

Examinar en ENASEM 2001--2021 la relación descriptiva entre expansión de
recepción pensionaria y apoyos públicos, brechas de cobertura y monto real,
incorporación y continuidad de receptores, identificando grupos sin vía propia
K/A y delimitando la sensibilidad de esa lectura a mortalidad registrada y
pérdida de observación.

\subsection{Objetivos específicos}

\begin{enumerate}
\item Distinguir pensión propia K, pensión del cónyuge, apoyo propio A y
prestación explícita de vejez A5, con sus superposiciones, negativas y
estados desconocidos.
\item Describir cobertura propia y diferencias por sexo, edad, escolaridad
inicial, tamaño de localidad y actividad, utilizando denominadores
compatibles con cada pregunta.
\item Comparar montos nominales y reales bajo una referencia de precios
común y valorar prestaciones parciales frente a una referencia económica
histórica compatible, sin inferir ingreso total o pobreza.
\item Contrastar nuevos receptores observados y continuos en las ventanas
monetarias centrales, examinando composición basal y evidencia institucional.
\item Distinguir variaciones de medianas transversales de cambios monetarios
individuales y evaluar su precisión y sensibilidad temporal.
\item Examinar quién permanece observable, fallece, responde por sustituto
o presenta información insuficiente, sin atribuir a los no observados
trayectorias monetarias desconocidas.
\end{enumerate}

\section{Hipótesis y proposiciones contrastables}
\label{sec:hipotesis}

La hipótesis general es que \emph{la ampliación de la cobertura no
necesariamente estuvo acompañada de una reducción equivalente de las
desigualdades monetarias ni de una mejora uniforme del poder adquisitivo}.
Se trata de una proposición a evaluar por vía, periodo y población, no de una
conclusión demostrada para todo 2001--2021.

Se organiza su evaluación en tres proposiciones:

\begin{enumerate}
\item[\textbf{H1.}] En ventanas de comparabilidad suficiente puede coexistir
mayor recepción con una variación no positiva de la mediana real entre
receptores. Se contrastan por separado variaciones de cobertura y de monto;
un contraejemplo con aumento real limita cualquier generalización.
\item[\textbf{H2.}] La distribución monetaria de los nuevos receptores puede
diferir de la de los continuos, y ese contraste puede variar entre K, A y
subgrupos institucionales. La diferencia mediana y sus intervalos se evalúan
sin presuponer que una coincidencia de medianas implique equivalencia.
\item[\textbf{H3.}] Una reducción de la brecha de cobertura entre hombres y
mujeres puede coexistir con diferencias monetarias persistentes entre
receptores. El cambio de la brecha de montos se estima directamente:
un intervalo que incluya cero no acredita un ensanchamiento o cierre.
\end{enumerate}

Estas proposiciones sistematizan una evidencia que ya fue explorada. No
corresponden a un protocolo preregistrado ni convierten los contrastes
seleccionados después del análisis en confirmación independiente. Los
intervalos son puntuales y condicionados a la muestra observada.
Las definiciones matemáticas y reglas para resultados nulos aparecen en el
capítulo~\ref{cap:datos}; el balance de evidencia y contraejemplos se desarrolla
en los capítulos~\ref{cap:resultados} y~\ref{cap:discusion}.

\section{Alcance y delimitaciones}

La población principal comprende personas registradas vivas, de 60 años o más
en cada visita, con entrevista directa. Se conserva el panel no balanceado:
cada corte incluye a quienes cumplen su elegibilidad y cada intervalo sigue
IDs verificables sin exigir seis presencias. Los sustitutos válidos se
mantienen como sensibilidad. El maestro completo de \cpTotalRegistros{}
personas-ronda y \cpTotalPersonas{} personas constituye el marco de integración,
no el denominador de todos los resultados.

K y A son conceptos operacionales de la encuesta. K no acredita
contributividad y A no equivale a una pensión no contributiva homogénea.
Los montos representan eventos captados o sumas personales parciales; no
ingreso económico completo. Se utilizan precios de septiembre de 2026 como
referencia común, sin incorporar una ronda adicional del estudio.

Las rupturas de cuestionario limitan la lectura temporal. K 2003--2012 y
K 2015--2018, así como A 2001--2003 y A 2018--2021, no se interpretan como
altas o bajas homologadas. La mortalidad y las pérdidas son sensibilidades
complementarias: no se estima supervivencia completa, ni se repondera por
observación ni se imputan recursos posteriores al fallecimiento.
Tampoco se desarrolla un modelo predictivo de vulnerabilidad.

\section{Contribución y organización de la tesis}

La literatura ya estudia seguridad económica, fuentes de ingreso,
dinámica transversal y longitudinal, y diferencias entre acceso y beneficio
pensionario \citep{degraff2018,perezarce2021,garciagarcia2025}. La aportación
específica se sitúa en el contraste conjunto de incorporación y continuidad,
montos reales, composición basal e instituciones para las dos ventanas
centrales, con incertidumbre que respeta las observaciones de una misma
persona. Es una ampliación empírica delimitada, no una pretensión de
prioridad conceptual frente a toda la literatura.

El capítulo~\ref{cap:marco} presenta conceptos, antecedentes y cronología
institucional. El capítulo~\ref{cap:datos} define fuentes, universos y
estimandos. El capítulo~\ref{cap:resultados} articula cobertura, montos,
incorporación y continuidad, preservando nulos y contraejemplos.
El capítulo~\ref{cap:discusion} discute el alcance de la hipótesis y de la
contribución. Los anexos conservan la descripción complementaria de CP1
y los pendientes reales de revisión.
